%I will structure the introduction as follows:

Ecological inference fallacy is the problem that arises when we are interested in relationships between variables at the individual level but we only observe data at the aggregate level.

A famous example of that problem is the one presented by \citet{Robinson1950} for the relationship between illiteracy and race, and illiteracy and migration in the United States. In that case, the author observed that areas with a higher proportion of Black population, as well as areas with a higher proportion of migrants, had lower illiteracy rates. However, at the individual level, Black people and migrants were more likely to be illiterate than White and non-migrant people. The explanation for this apparent paradox is that Black and migrant people were more likely to live in or move to areas with lower illiteracy rates, and therefore the aggregate data masked the individual-level relationship.

We can also think of economic examples.

Consider, for instance, a retail store chain that wants to investigate the consumption habits of young people by observing sales data from different stores.
The retail store chain might know the quantities of goods sold in different stores and, at the same time, demographic data about consumers in the different areas served by those stores.
Therefore, it would be tempting to link the percentage of young people to the quantities of different goods sold, suggesting that goods that sell more in stores with younger consumers are purchased by young consumers.

However, that could lead to false findings. In fact, when we consider data on cities, it might well be that stores with a greater share of young consumers are located in central areas near universities, where richer people also usually live. Therefore, even if young people might be more price-sensitive consumers, their habits would be masked by their neighbors' habits.
% (often times, these problems are overcome by including covariates, fixed effects or instrumental variables, but these solutions are not always apt to solve problems of ecological inference, where the problem is that we do not observe individual-level data, and therefore we cannot control for individual-level confounders)

In any case, the most obvious application of ecological inference concerns electoral behavior. We are interested in the relationship between individual-level characteristics and voting behavior, but, obviously (and fortunately!), we cannot observe how individuals vote. We are therefore unable to link voting behavior directly to individual-level characteristics.
In these cases, we are forced to rely on aggregate data, ideally at the most granular level of aggregation available.
However, even at the most granular level of aggregation, individual-level relationships may be severely masked.

This issue is particularly relevant for political economists, who often study the relationships between voters' characteristics, electoral outcomes, and public policies.

\citet{kuriwaki2026roleconfounderslinearityecological} present a particularly convincing example of this problem using the 1968 US presidential election. In addition to the Republican and Democratic candidates, Richard Nixon and Hubert Humphrey, respectively, there was a relevant third-party candidate, George Wallace, who led the segregationist American Independent Party.
George Wallace performed best in the Southern states, where the proportion of the Black population was generally higher than in the rest of the country. Therefore, a naive analysis of the aggregate data might suggest that Black people were more likely to vote for George Wallace, since states with larger Black populations also had higher proportions of votes for him.
That would, of course, be a false conclusion: Black people certainly did not favor George Wallace, whose platform was segregationist and overtly racist.

\citet{kuriwaki2026roleconfounderslinearityecological} go on to show that, even at more granular levels of aggregation, such as the county level, the relationship between the proportion of the Black population and the proportion of votes for George Wallace remains positive when considering all counties in the nation. Moreover, even when restricting the analysis to specific states, such as North Carolina, we still observe a positive relationship: the higher the proportion of the Black population in a county, the higher the proportion of votes for George Wallace.

The key is that, while Black voters were relatively homogeneous in their voting behavior—generally favoring the Democratic candidate and certainly not voting for George Wallace—White voters were much more heterogeneous. In particular, White people with segregationist views were more likely to live in areas with a higher proportion of Black residents.

The same mechanism is at work in the retail-store example: while young consumers might be relatively homogeneous in their spending habits, the larger group of older consumers may have different habits and drive the aggregate relationship.


\citet{Robinson1950}, who raised the issue of ecological fallacy with the illiteracy example, came to the conclusion that ecological fallacy would constitute an inextricable problem and he concluded his paper telling that he was 
\begin{quotation} ... aware that this conclusion has serious consequences, and that its effect appears wholly negative because it throws serious doubt upon the validity of a number of important studies made in recent years. The purpose of this paper will have been accomplished, however, if it prevents the future computation of meaningless correlations and stimulates the study of similar problems with the use of meaningful correlations between the properties of individuals.
\end{quotation}

Econometrics in 1950, however, had considerably fewer tools than it does today, and we might therefore think that some of these techniques could be successfully deployed to address the ecological fallacy.

In particular, we might be tempted to argue that the problem could be solved by including covariates.

That, however, is not necessarily sufficient.
In fact, when \citet{kuriwaki2026roleconfounderslinearityecological} examine the results of presidential elections in North Carolina, they are implicitly including a state indicator and interacting it with the proportion of Black voters.
Indeed, their estimate can be interpreted as the coefficient on the share of Black voters, conditional on the county being located in North Carolina.

Analogously, we might include every possible interaction between variables and functional form, or impose a nonlinear relationship between the share of Black voters and electoral outcomes. This could produce a positive relationship between the share of Black voters and segregationist votes at relatively low levels of Black population, which would reverse once Black voters became the majority.

Another method we might consider is the inclusion of fixed effects.
However, fixed effects require repeated observations of the same unit, which are not available in our examples.

All these examples highlight the common problem with using traditional econometric techniques to address ecological inference: the discrepancy between the unit of observation and the quantity we seek to estimate.

Indeed, each of these techniques may help identify relationships between variables at the level of the unit of observation, but none would solve the ecological fallacy of composition.

Even our most naive estimates may correctly describe relationships at the level of the units of observation.
It may be true that counties with more migrants have lower illiteracy rates, even after accounting for other variables or effects; that consumers in areas with younger populations spend more on expensive goods; or that precincts with higher shares of Black voters favor segregationist candidates.

What these estimates cannot tell us is whether migrants have lower or higher illiteracy rates, whether younger consumers spend more or less on expensive goods, or how Black people vote.

To answer these questions, we must introduce notation that explicitly accounts for the aggregate nature of our units of observation and impose a set of assumptions.

Only then can we move from the agnostic conclusion of \citet{Robinson1950} to ecological estimates.

In the next section, I will introduce a general notation for ecological inference problems and show how different assumptions about the distribution of individual-level characteristics within geographical units of observation can lead to different ecological inference methods.

In the following section, I will present a more specific model for electoral analysis.

Both sections draw heavily on \citet{kuriwaki2026roleconfounderslinearityecological}, which reviews ecological inference methods, and on \citet{mccartan2025identificationsemiparametricestimationconditional}, which presents the semiparametric method for ecological inference implemented in the R package \texttt{seine} (\citet{seine}) and used in this dissertation.

The notation I present is nevertheless slightly different and better suited to describing the issues considered in this dissertation.

After presenting the model, I will use synthetic examples to illustrate how the method works and assess its performance in a controlled environment.

Finally, I will apply the method to real data from the 2020 US presidential election in Texas.

%In particular, the inclusion of covariates or fixed effects or instrumental variables

%These techniques in any case can possibly work well when we apply them to individual-level data, while they are not apt to confront aggregate data.
%Think again of the previous examples (eh ci devo pensare anche io)

%I might also consider BLP

%It is therefore useful to introduce a new notation for the problem of composition






%Introduction: \citet{kuriwaki2026roleconfounderslinearityecological}, but also \citet{mccartan2025identificationsemiparametricestimationconditional}.


%I should say something about the more general importance of estimating correctly when we do not observe individual-level data.

%I should say that sometimes we do not observe individual-level data and we might be tempted to infere something based on highly granular data. But these inferences may be wrong.

%That could be relevant for economists in various applications and analyses and it is not solely a nuisance for econometricians.

%In fact, data are usually reported by official bureau of statistics even at highly granular level but not at individual level and some relationships become masked.


%Subsequently, I should present the problem with appropriate notation.

%And probably I should also mark the difference between this problem and others problem of aggregation such as the one discussed in BLP model.

%In any case, I will later decide how to structure the dissertation and where to put the complete explanation of the method: in principle, I would put it into the introduction and later I would present the syntethic example just referring to what I have already written here.


%In particolare, quali sono le differenze fra questo metodo e gli altri?

%Rispetto agli effetti fissi, qua dico che la variabilità non è una semplice trasposizione data dagli effetti fissi (dico ad esempio che tra Pennsylvania e New Mexico in New Mexico i laureati votano il 10\% più Democratico, ma non perché il New Mexico sia spostato tutto 10\% più Democratico; magari lo sono i laureati, mentre i non laureati sono il 5\% meno Democratici rispetto a quelli della Pennsylvania).

%Invece Gary King proponeva di fare i coefficienti Beta con una certa variabilità nelle diverse geografie ma secondo una distribuzione parametrica (Normale).

%E così via, mostrando un po' di quello che mostrano in \citet{kuriwaki2026roleconfounderslinearityecological}.

%Un'altra possibilità, direi, è quella di fare una regressione usando tanti regressori, dicendo che quello precedente non è un problema di ecological fallacy ma di omitted variables.

%Posso dire che l'altra differenza è quella con l'approccio di modellare le interazioni tra agenti che è molto più complicato.
